\section{Values in the Extended Reals}
\label{app:ennreal}

Valuations take values in the extended non-negative reals $\ENNReal$,
following the measure-theoretic convention of \textsf{mathlib}.  
As a result,
sums indexed by an arbitrary family are unconditionally defined, with no
summability hypothesis, and they commute over dependent sums: for a
function $f$ on a sigma type $\Sigma\,a,\,\beta(a)$, one has
$\sum_{p} f(p) = \sum_{a} \sum_{b} f(a, b)$ with no side condition.  This
is what lets the countable-mixture results of
Section~\ref{sec:monad} hold as stated.  Additionally, subtraction is
truncated at $0$.  This never affects the main text, since every
subtraction there is applied to a pair already ordered by monotonicity,
as with the mass of Section~\ref{sec:representation}, so the truncation
is never triggered.

\section{Lean Statements}
\label{app:lean}

\subsection{Introduction}

\subsection{Valuations on Frames}

\phantomsection\label{lean:valuation}
\begin{lstlisting}
structure Valuation (Ω : Type*) [Order.Frame Ω] where
  toFun : Ω → ℝ≥0∞
  map_bot' : toFun ⊥ = 0
  map_top' : toFun ⊤ = 1
  mono' : Monotone toFun
  modular' : ∀ a b, toFun a + toFun b = toFun (a ⊔ b) + toFun (a ⊓ b)
\end{lstlisting}

\phantomsection\label{lean:add-compl-le-one}
\begin{lstlisting}
theorem Valuation.add_compl_le_one (v : Valuation Ω) (a : Ω) :
    v a + v aᶜ ≤ 1
\end{lstlisting}

\phantomsection\label{lean:valuation-add-compl-eq-sup}
\begin{lstlisting}
theorem Valuation.add_compl_eq_sup (v : Valuation Ω) (a : Ω) :
    v a + v aᶜ = v (a ⊔ aᶜ)
\end{lstlisting}

\phantomsection\label{lean:valuation-le-one}
\begin{lstlisting}
theorem Valuation.le_one (v : Valuation Ω) (a : Ω) : v a ≤ 1
\end{lstlisting}

\phantomsection\label{lean:valuation-slack-eq-dngap-add-demorgangap}
\begin{lstlisting}
theorem Valuation.slack_eq_dnGap_add_deMorganGap (v : Valuation Ω) (a : Ω) :
    v.slack a = v.dnGap a + v.deMorganGap a
\end{lstlisting}

\phantomsection\label{lean:valuation-slack-eq-sub-add-sub}
\begin{lstlisting}
theorem Valuation.slack_eq_sub_add_sub (v : Valuation Ω) (a : Ω) :
    v.slack a = (v aᶜᶜ - v a) + (1 - (v aᶜᶜ + v aᶜ))

theorem Valuation.le_compl_compl (v : Valuation Ω) (a : Ω) : v a ≤ v aᶜᶜ
\end{lstlisting}

\phantomsection\label{lean:valuation-slack-submodular}
\begin{lstlisting}
theorem Valuation.slack_submodular (v : Valuation Ω) (a b : Ω) :
    v.slack (a ⊔ b) + v.slack (a ⊓ b) ≤ v.slack a + v.slack b
\end{lstlisting}

\phantomsection\label{lean:isdemorgan}
\begin{lstlisting}
def Order.Frame.IsDeMorgan (Ω : Type*) [Order.Frame Ω] : Prop :=
  ∀ a : Ω, aᶜᶜ ⊔ aᶜ = ⊤
\end{lstlisting}

\phantomsection\label{lean:valuation-slack-eq-dngap-of-isdemorgan}
\begin{lstlisting}
theorem Valuation.slack_eq_dnGap_of_isDeMorgan
    (hΩ : Order.Frame.IsDeMorgan Ω) (v : Valuation Ω) (a : Ω) :
    v.slack a = v.dnGap a
\end{lstlisting}

\phantomsection\label{lean:condval-mul}
\begin{lstlisting}
noncomputable def Valuation.condVal (v : Valuation Ω) (b : Ω) (hb : v b ≠ 0) :
    Valuation Ω

theorem Valuation.condVal_mul (v : Valuation Ω) (b : Ω) (hb : v b ≠ 0) (a : Ω) :
    v.condVal b hb a * v b = v (a ⊓ b)

theorem Valuation.condVal_symm (v : Valuation Ω) {a b : Ω}
    (ha : v a ≠ 0) (hb : v b ≠ 0) :
    v.condVal b hb a * v b = v.condVal a ha b * v a
\end{lstlisting}

\phantomsection\label{lean:cond-add-compl-le}
\begin{lstlisting}
theorem Valuation.cond_add_compl_le (v : Valuation Ω) (a b : Ω) :
    v (a ⊓ b) + v (aᶜ ⊓ b) ≤ v b
\end{lstlisting}

\phantomsection\label{lean:add-compl-eq-one-of-complemented}
\begin{lstlisting}
theorem Valuation.add_compl_eq_one_of_complemented (v : Valuation Ω) {a : Ω}
    (ha : a ⊔ aᶜ = ⊤) : v a + v aᶜ = 1
\end{lstlisting}

\phantomsection\label{lean:classical-additivity}
\begin{lstlisting}
variable [CompleteBooleanAlgebra Ω]

theorem classical_additivity (v : Valuation Ω) (a : Ω) :
    v a + v aᶜ = 1
\end{lstlisting}

\phantomsection\label{lean:classical-slack-zero}
\begin{lstlisting}
variable [CompleteBooleanAlgebra Ω]

theorem classical_slack_zero (v : Valuation Ω) (a : Ω) :
    v.slack a = 0
\end{lstlisting}

\phantomsection\label{lean:valuation-mix}
\begin{lstlisting}
noncomputable def Valuation.mix (w : ι → ℝ≥0∞) (hw : ∑ i, w i = 1)
    (v : ι → Valuation Ω) : Valuation Ω
\end{lstlisting}

\subsection{Where Classicality Lives}

\phantomsection\label{lean:hinge}
\begin{lstlisting}
def Valuation.HasClassicalNegation (v : Valuation Ω) : Prop :=
  ∀ a, v a + v aᶜ = 1

theorem hasClassicalNegation_of_em (hem : ∀ a : Ω, a ⊔ aᶜ = ⊤)
    (v : Valuation Ω) : v.HasClassicalNegation

theorem em_of_forall_hasClassicalNegation
    (h : ∀ v : Valuation Ω, v.HasClassicalNegation) : ∀ a : Ω, a ⊔ aᶜ = ⊤
\end{lstlisting}

\phantomsection\label{lean:sharp-iff-point}
\begin{lstlisting}
theorem sharp_iff_point (v : Valuation Ω) :
    v.IsSharp ↔ ∃ (J : Order.Ideal Ω) (hJ : J.IsPrime) (htop : ⊤ ∉ J),
      v = Ideal.toValuation J hJ htop
\end{lstlisting}

\phantomsection\label{lean:sharp-scott-iff-completelyprimepoint}
\begin{lstlisting}
theorem sharp_scott_iff_completelyPrimePoint (v : Valuation Ω) :
    v.IsSharp ∧ v.ScottContinuous ↔
      ∃ (J : Order.Ideal Ω) (hJ : Ideal.IsCompletelyPrime J) (htop : ⊤ ∉ J),
        v = Ideal.toValuation J hJ.1 htop
\end{lstlisting}

\subsection{Valuation Representation}

\phantomsection\label{lean:valuation-mass}
\begin{lstlisting}
noncomputable def Valuation.mass (v : Valuation (LowerSet P)) (p : P) : ℝ≥0∞
\end{lstlisting}

\phantomsection\label{lean:valuation-eq-sum-mass}
\begin{lstlisting}
theorem Valuation.eq_sum_mass (v : Valuation (LowerSet P)) (U : LowerSet P) :
    v U = ∑ p ∈ U.toFinset, v.mass p

theorem Valuation.sum_mass (v : Valuation (LowerSet P)) : ∑ p, v.mass p = 1
\end{lstlisting}

\phantomsection\label{lean:valuation-topmf}
\begin{lstlisting}
noncomputable def Valuation.toPMF (v : Valuation (LowerSet P)) : PMF P
\end{lstlisting}

\phantomsection\label{lean:valuation-exists-pmf-presentation}
\begin{lstlisting}
theorem Valuation.exists_pmf_presentation {Θ : Type u} [Order.Frame Θ] [Finite Θ]
    (v : Valuation Θ) :
    ∃ (P : Type u) (_ : PartialOrder P) (_ : Fintype P) (_ : DecidableEq P)
      (e : Θ ≃o LowerSet P),
      (∀ a, v a = ∑ p ∈ (e a).toFinset, (v.mapEquiv e).mass p) ∧
        ∑ p, (v.mapEquiv e).mass p = 1
\end{lstlisting}

\phantomsection\label{lean:valuation-eq-mix-deltapoint}
\begin{lstlisting}
noncomputable def deltaPoint (p : P) : Valuation (LowerSet P)

theorem Valuation.eq_mix_deltaPoint (v : Valuation (LowerSet P)) :
    v = Valuation.mix v.mass v.sum_mass (fun p => deltaPoint p)
\end{lstlisting}

\phantomsection\label{lean:exists-valuation-not-point-representable}
\begin{lstlisting}
theorem exists_valuation_not_point_representable :
    ∃ v : Valuation (LowerSet ℕ), ∑' p, v.mass p ≠ v ⊤

theorem topIndicator_mass (n : ℕ) : topIndicator.mass n = 0
\end{lstlisting}

\phantomsection\label{lean:topindicator-not-scott}
\begin{lstlisting}
theorem topIndicator_not_scott :
    (⨆ n, topIndicator (LowerSet.Iic n)) ≠ topIndicator ⊤
\end{lstlisting}

\phantomsection\label{lean:eq-tsum-mass-of-scott}
\begin{lstlisting}
theorem eq_tsum_mass_of_scott (v : Valuation (LowerSet ℕ))
    (hsc : v ⊤ = ⨆ n, v (LowerSet.Iic n)) : v ⊤ = ∑' n, v.mass n

variable [∀ p : P, Finite (Set.Iic p)]

theorem Valuation.isPurelyAtomic_of_scott (v : Valuation (LowerSet P))
    (hsc : v ⊤ = ⨆ (S : Finset P), v (S.sup LowerSet.Iic)) : IsPurelyAtomic v
\end{lstlisting}

\phantomsection\label{lean:tsum-mass-le}
\begin{lstlisting}
theorem tsum_mass_le (v : Valuation (LowerSet P)) : ∑' p, v.mass p ≤ v ⊤
\end{lstlisting}

\subsection{The Dempster--Shafer Reading}

\phantomsection\label{lean:inclusion-exclusion}
\begin{lstlisting}
theorem Valuation.inclusion_exclusion (v : Valuation Ω) (s : Finset ι)
    (a : ι → Ω) : v (s.sup a) + v.iesEven s a = v.iesOdd s a
\end{lstlisting}

\phantomsection\label{lean:infty-monotone}
\begin{lstlisting}
theorem Valuation.infty_monotone (v : Valuation Ω) (s : Finset ι) (a : ι → Ω) :
    v.iesOdd s a ≤ v ((s.sup a)ᶜᶜ) + v.iesEven s a
\end{lstlisting}

\phantomsection\label{lean:two-monotone}
\begin{lstlisting}
theorem Valuation.two_monotone (v : Valuation Ω) (a b : Ω) :
    v a + v b ≤ v (a ⊔ b)ᶜᶜ + v (a ⊓ b)
\end{lstlisting}

\phantomsection\label{lean:plausibility-sub-self}
\begin{lstlisting}
noncomputable def Valuation.plausibility (v : Valuation Ω) (a : Ω) : ℝ≥0∞ := 1 - v aᶜ

theorem Valuation.plausibility_sub_self (v : Valuation Ω) (a : Ω) :
    v.plausibility a - v a = v.slack a
\end{lstlisting}

\phantomsection\label{lean:sup-compl-decomp}
\begin{lstlisting}
theorem Valuation.sup_compl_decomp (v : Valuation Ω) (a b : Ω) :
    v (a ⊔ bᶜ) = v bᶜ + (v (a ⊓ b) + v.emGap a b)
\end{lstlisting}

\phantomsection\label{lean:dempstercond}
\begin{lstlisting}
noncomputable def Valuation.dempsterCond (v : Valuation Ω) (b : Ω)
    (hb : v b ≠ 0) : Valuation Ω

theorem Valuation.dempsterCond_apply (v : Valuation Ω) (b : Ω) (hb : v b ≠ 0)
    (a : Ω) : v.dempsterCond b hb a = (v (a ⊔ bᶜ) - v bᶜ) / (1 - v bᶜ)
\end{lstlisting}

\phantomsection\label{lean:dempstercond-eq-condval-iff-slack}
\begin{lstlisting}
theorem Valuation.dempsterCond_eq_condVal_iff_slack (v : Valuation Ω) {b : Ω}
    (hb : v b ≠ 0) :
    (∀ a, v.dempsterCond b hb a = v.condVal b hb a) ↔ v.slack b = 0
\end{lstlisting}

\subsection{Cox's Theorem}

\phantomsection\label{lean:coxmodel}
\begin{lstlisting}
structure CoxModel (Ω : Type*) [Order.Frame Ω] where
  pl : Ω → Ω → ℝ
  mono_left : ∀ c, Monotone fun a => pl a c
  pl_bot : ∀ c, pl ⊥ c = 0
  pl_top : ∀ c, pl ⊤ c = 1
  F : ℝ → ℝ → ℝ
  conj : ∀ a b c, pl (a ⊓ b) c = F (pl a (b ⊓ c)) (pl b c)
  F_assoc : ∀ x y z, F (F x y) z = F x (F y z)
  F_strictMono_left : ∀ y, 0 < y → StrictMono fun x => F x y
\end{lstlisting}

\phantomsection\label{lean:isdemorgan-lowersetv}
\begin{lstlisting}
theorem isDeMorgan_lowerSetV : Order.Frame.IsDeMorgan (LowerSet V)
\end{lstlisting}

\phantomsection\label{lean:no-disjunction-functional}
\begin{lstlisting}
theorem modularity_irreducible :
    ∃ (Ω : Type) (_ : Order.Frame Ω) (q : Ω → ℝ≥0∞),
      Monotone q ∧ q ⊥ = 0 ∧ q ⊤ = 1 ∧
      (∀ x y : Ω, x ⊓ y = ⊥ → q (x ⊔ y) = q x + q y) ∧
      ¬ (∀ x y : Ω, q x + q y = q (x ⊔ y) + q (x ⊓ y))

theorem no_disjunction_functional :
    ¬ ∃ S : ℝ≥0∞ → ℝ≥0∞ → ℝ≥0∞,
        ∀ x y : LowerSet V, qV (x ⊔ y) = S (qV x) (qV y)
\end{lstlisting}

\phantomsection\label{lean:constructive-cox}
\begin{lstlisting}
theorem constructive_cox (M : CoxModel Ω)
    (hmod : ∀ a b : Ω, M.pl a ⊤ + M.pl b ⊤ = M.pl (a ⊔ b) ⊤ + M.pl (a ⊓ b) ⊤) :
    ∃ (v : Valuation Ω) (g : ℝ → ℝ≥0∞),
      StrictMonoOn g (Set.Icc 0 1) ∧ g 0 = 0 ∧ g 1 = 1 ∧
        ∀ a, v a = g (M.pl a ⊤)
\end{lstlisting}

\phantomsection\label{lean:coxmodelennreal}
\begin{lstlisting}
noncomputable def coxModelENNReal : CoxModel ℝ≥0∞
\end{lstlisting}

\phantomsection\label{lean:exists-mul-generator}
\begin{lstlisting}
theorem Scale.exists_mul_generator :
    ∃ G : ℝ → ℝ,
      (∀ x y, S.u ≤ x → S.u ≤ y → G (S.F x y) = G x * G y) ∧
      (∀ x y, S.u ≤ x → x < y → G x < G y)

theorem Scale.aczelStatement_cone :
    ∃ G : ℝ → ℝ,
      (∀ x y, S.u ≤ x → S.u ≤ y → G (S.F x y) = G x * G y) ∧
      StrictMonoOn G (Set.Ici S.u)
\end{lstlisting}

\phantomsection\label{lean:exists-diag-eq}
\begin{lstlisting}
theorem exists_diag_eq {F : ℝ → ℝ → ℝ} {a b c : ℝ} (hab : a ≤ b)
    (hcont : ContinuousOn (fun x => F x x) (Set.Icc a b))
    (hc : c ∈ Set.Icc (F a a) (F b b)) : ∃ x ∈ Set.Icc a b, F x x = c
\end{lstlisting}

\phantomsection\label{lean:scale}
\begin{lstlisting}
structure Scale where
  F : ℝ → ℝ → ℝ
  hdiv : ∀ c, ∃ x, F x x = c
  hassoc : ∀ x y z, F (F x y) z = F x (F y z)
  hpos : ∀ c x, x < F x c
  hmono : ∀ x₁ x₂ y₁ y₂ : ℝ, x₁ ≤ x₂ → y₁ ≤ y₂ → F x₁ y₁ ≤ F x₂ y₂
  hcont : Continuous fun p : ℝ × ℝ => F p.1 p.2
  hbot : ∀ x, Filter.Tendsto (fun z => F x z) Filter.atBot (nhds x)
  u : ℝ
\end{lstlisting}

\subsection{Two Grounding Models}

\phantomsection\label{lean:measure-tovaluationopens}
\begin{lstlisting}
noncomputable def MeasureTheory.Measure.toValuationOpens
    (μ : Measure X) [IsProbabilityMeasure μ] : Valuation (Opens X)
\end{lstlisting}

\phantomsection\label{lean:tovaluationopens-slack-eq-frontier}
\begin{lstlisting}
theorem toValuationOpens_slack_eq_frontier (μ : Measure X)
    [IsProbabilityMeasure μ] (U : Opens X) :
    μ.toValuationOpens.slack U = μ (frontier (U : Set X))
\end{lstlisting}

\phantomsection\label{lean:classical-posterior-compl-eq}
\begin{lstlisting}
theorem classical_posterior_compl_eq (μ : Measure X)
    [IsProbabilityMeasure μ] (a b : Opens X)
    (hb : μ.toValuationOpens b ≠ 0) :
    μ ((a : Set X)ᶜ ∩ b) / μ b
      = μ.toValuationOpens.condVal b hb aᶜ
        + μ (frontier (a : Set X) ∩ b) / μ b
\end{lstlisting}

\phantomsection\label{lean:interiormeasure-add-compl-le}
\begin{lstlisting}
theorem interiorMeasure_add_compl_le (μ : Measure X) [IsProbabilityMeasure μ]
    (S : Set X) : μ.interiorMeasure S + μ.interiorMeasure Sᶜ ≤ 1
\end{lstlisting}

\phantomsection\label{lean:isdemorgan-opens-iff-extremallydisconnected}
\begin{lstlisting}
theorem isDeMorgan_opens_iff_extremallyDisconnected :
    Order.Frame.IsDeMorgan (Opens X) ↔ ExtremallyDisconnected X
\end{lstlisting}

\phantomsection\label{lean:isdemorgan-opens-iff-stonean}
\begin{lstlisting}
theorem isDeMorgan_opens_iff_stonean [CompactSpace X] [T2Space X] :
    Order.Frame.IsDeMorgan (Opens X) ↔ ExtremallyDisconnected X
\end{lstlisting}

\phantomsection\label{lean:frontier-subset-ptz}
\begin{lstlisting}
inductive ThreePt : Type
  | ptA : ThreePt
  | ptB : ThreePt
  | ptZ : ThreePt

theorem frontier_subset_ptZ {U : Set ThreePt} (hU : IsOpen U) :
    frontier U ⊆ {ptZ}
\end{lstlisting}

\phantomsection\label{lean:exists-ne-valuation-same-slack}
\begin{lstlisting}
theorem exists_ne_valuation_same_slack :
    ∃ v v' : Valuation (Opens ThreePt), v.slack = v'.slack ∧ v ≠ v'
\end{lstlisting}

\phantomsection\label{lean:haltsopen-compl}
\begin{lstlisting}
theorem haltsOpen_compl : haltsOpenᶜ = (⊥ : Opens Prop)
\end{lstlisting}

\phantomsection\label{lean:slack-eq-boundary}
\begin{lstlisting}
theorem slack_eq_boundary {p : ℝ≥0∞} (hp1 : p ≤ 1) :
    (sierpinskiValuation hp1).slack haltsOpen
      = coinMeasure p (frontier {q : Prop | q})
\end{lstlisting}

\phantomsection\label{lean:sharpreadout-not-computable}
\begin{lstlisting}
noncomputable def haltingWeight (n : ℕ) (c : Code) : ℝ≥0∞ :=
  if (c.eval n).Dom then 1 else 0

noncomputable def sharpReadout (n : ℕ) (c : Code) : Valuation (Opens Prop) :=
  sierpinskiValuation (haltingWeight_le_one n c)

theorem sharpReadout_not_computable (n : ℕ) :
    ¬ComputablePred fun c : Code => sharpReadout n c haltsOpen = 1
\end{lstlisting}

\phantomsection\label{lean:isopen-wordcylinder}
\begin{lstlisting}
theorem isOpen_wordCylinder (w : List Bool) : IsOpen (wordCylinder w)

theorem cantorMeasure_wordCylinder (w : List Bool) :
    cantorMeasure (wordCylinder w) = (1 / 2) ^ w.length
\end{lstlisting}

\phantomsection\label{lean:cantormeasure-biunion-wordcylinder}
\begin{lstlisting}
theorem disjoint_wordCylinder_of_not_prefix {v w : List Bool}
    (hv : ¬ v <+: w) (hw : ¬ w <+: v) :
    Disjoint (wordCylinder v) (wordCylinder w)

theorem cantorMeasure_biUnion_wordCylinder {F : Finset (List Bool)}
    (hF : PrefixFree F) :
    cantorMeasure (⋃ w ∈ F, wordCylinder w) = ∑ w ∈ F, (1 / 2 : ℝ≥0∞) ^ w.length
\end{lstlisting}

\subsection{The Monad Question}

\phantomsection\label{lean:valuation-comap}
\begin{lstlisting}
def Valuation.comap (f : FrameHom Γ Ω) (v : Valuation Ω) : Valuation Γ

theorem Valuation.comap_id (v : Valuation Ω) : v.comap (FrameHom.id Ω) = v

theorem Valuation.comap_comp (f : FrameHom Γ Ω) (g : FrameHom Δ Γ)
    (v : Valuation Ω) : v.comap (f.comp g) = (v.comap f).comap g
\end{lstlisting}

\phantomsection\label{lean:exists-rect-decomposition}
\begin{lstlisting}
theorem exists_rect_decomposition (W : LowerSet (ℕ × ℕ)) :
    ∃ (N : ℕ) (A B : ℕ → LowerSet ℕ),
      W = (Finset.range N).sup fun i => A i ×ˢ B i
\end{lstlisting}

\phantomsection\label{lean:eq-on-finsetsup-of-eq-on}
\begin{lstlisting}
theorem Valuation.eq_on_finsetSup_of_eq_on {G : Set Ω}
    (hG : ∀ g ∈ G, ∀ g' ∈ G, g ⊓ g' ∈ G) (u u' : Valuation Ω)
    (h : ∀ g ∈ G, u g = u' g) {ι : Type*} [DecidableEq ι]
    (s : Finset ι) (a : ι → Ω) (ha : ∀ i ∈ s, a i ∈ G) :
    u (s.sup a) = u' (s.sup a)
\end{lstlisting}

\phantomsection\label{lean:valuation-product-unique}
\begin{lstlisting}
theorem Valuation.product_unique (v w : Valuation (LowerSet ℕ))
    {u u' : Valuation (LowerSet (ℕ × ℕ))}
    (hu : ∀ A B : LowerSet ℕ, u (A ×ˢ B) = v A * w B)
    (hu' : ∀ A B : LowerSet ℕ, u' (A ×ˢ B) = v A * w B) : u = u'
\end{lstlisting}

\phantomsection\label{lean:ultrafiltervaluation}
\begin{lstlisting}
noncomputable def ultrafilterValuation (W : Ultrafilter ι) : Valuation (Set ι)
\end{lstlisting}

\phantomsection\label{lean:rect-mem-iff}
\begin{lstlisting}
theorem rect_mem_iff {α β : Type*} {W : Ultrafilter (α × β)}
    {S : Set α} {T : Set β} :
    S ×ˢ T ∈ W ↔ S ∈ W.map Prod.fst ∧ T ∈ W.map Prod.snd
\end{lstlisting}

\phantomsection\label{lean:product-valuation-not-unique}
\begin{lstlisting}
theorem product_valuation_not_unique :
    ∃ u u' : Valuation (Set (ℕ × ℕ)),
      (∀ S T : Set ℕ, u (S ×ˢ T) = u' (S ×ˢ T)) ∧ u ≠ u'
\end{lstlisting}

\phantomsection\label{lean:fubini-fails-for-valuations}
\begin{lstlisting}
theorem fubini_fails_for_valuations :
    ∃ (v w : Valuation (Set ℕ)) (u u' : Valuation (Set (ℕ × ℕ))),
      (∀ S T : Set ℕ, u (S ×ˢ T) = v S * w T) ∧
      (∀ S T : Set ℕ, u' (S ×ˢ T) = v S * w T) ∧ u ≠ u'
\end{lstlisting}

\phantomsection\label{lean:mixcountable}
\begin{lstlisting}
noncomputable def Valuation.mixCountable {ι : Type*} (w : ι → ℝ≥0∞)
    (hw : ∑' i, w i = 1) (v : ι → Valuation Ω) : Valuation Ω
\end{lstlisting}

\phantomsection\label{lean:mixcountable-mixcountable}
\begin{lstlisting}
theorem Valuation.mixCountable_mixCountable {ι : Type*} {κ : ι → Type*}
    (w : ι → ℝ≥0∞) (hw : ∑' i, w i = 1)
    (u : ∀ i, κ i → ℝ≥0∞) (hu : ∀ i, ∑' j, u i j = 1)
    (v : ∀ i, κ i → Valuation Ω) :
    Valuation.mixCountable w hw
        (fun i => Valuation.mixCountable (u i) (hu i) (v i))
      = Valuation.mixCountable (fun p : Σ i, κ i => w p.1 * u p.1 p.2)
          (tsum_sigma_mul_eq_one w hw u hu) (fun p => v p.1 p.2)
\end{lstlisting}

\phantomsection\label{lean:mixcountable-unique}
\begin{lstlisting}
theorem Valuation.mixCountable_unique (v : Valuation Ω) :
    Valuation.mixCountable (fun _ : PUnit => 1) (by simp) (fun _ => v) = v
\end{lstlisting}

\phantomsection\label{lean:mixcountable-comm}
\begin{lstlisting}
theorem Valuation.mixCountable_comm {ι κ : Type*}
    (w : ι → ℝ≥0∞) (hw : ∑' i, w i = 1) (u : κ → ℝ≥0∞) (hu : ∑' j, u j = 1)
    (v : ι → κ → Valuation Ω) :
    Valuation.mixCountable w hw (fun i => Valuation.mixCountable u hu (v i))
      = Valuation.mixCountable u hu
          (fun j => Valuation.mixCountable w hw (fun i => v i j))
\end{lstlisting}

\subsection{Valuations Beyond Frames}

\noindent Theorem~\ref{thm:sharp-separating}.
\phantomsection\label{lean:sharp-separating}
\begin{lstlisting}
theorem exists_sharp_separating {a b : Ω} (hab : ¬ a ≤ b) :
    ∃ v : Valuation Ω, v.IsSharp ∧ v a = 1 ∧ v b = 0
\end{lstlisting}

\noindent Definition~\ref{def:lvaluation}.
\phantomsection\label{lean:lvaluation}
\begin{lstlisting}
structure LValuation (Ω : Type*) [Lattice Ω] [BoundedOrder Ω] where
  toFun : Ω → ℝ≥0∞
  map_bot' : toFun ⊥ = 0
  map_top' : toFun ⊤ = 1
  mono'    : Monotone toFun
  modular' : ∀ a b, toFun a + toFun b
             = toFun (a ⊔ b) + toFun (a ⊓ b)
\end{lstlisting}

\noindent Proposition~\ref{prop:lvaluation-comap}.
\phantomsection\label{lean:lvaluation-comap}
\begin{lstlisting}
def LValuation.comap (f : BoundedLatticeHom Γ Ω) (v : LValuation Ω) :
    LValuation Γ

theorem LValuation.comap_id (v : LValuation Ω) :
    v.comap (BoundedLatticeHom.id Ω) = v

theorem LValuation.comap_comp (f : BoundedLatticeHom Γ Ω)
    (g : BoundedLatticeHom Δ Γ) (v : LValuation Ω) :
    v.comap (f.comp g) = (v.comap f).comap g
\end{lstlisting}

\noindent Theorem~\ref{thm:lvaluation-separating}.
\phantomsection\label{lean:lvaluation-separating}
\begin{lstlisting}
theorem LValuation.exists_sharp_separating [DistribLattice Ω]
    {a b : Ω} (hab : ¬ a ≤ b) :
    ∃ v : LValuation Ω, v.IsSharp ∧ v a = 1 ∧ v b = 0
\end{lstlisting}

\noindent Definition~\ref{def:m3}.
\phantomsection\label{lean:m3}
\begin{lstlisting}
inductive M3 : Type
  | z | a | b | c | o

instance : Lattice M3 where
  -- z bottom, o top, a b c pairwise-incomparable atoms;
  -- any two distinct atoms join to o and meet to z

instance : BoundedOrder M3 where
  top := o
  bot := z

theorem M3.not_distributive :
    ¬ ∀ x y w : M3, x ⊓ (y ⊔ w) = (x ⊓ y) ⊔ (x ⊓ w)
\end{lstlisting}

\noindent Theorem~\ref{thm:m3-no-sharp}.
\phantomsection\label{lean:m3-no-sharp}
\begin{lstlisting}
theorem M3.no_sharp_valuation (v : LValuation M3) : ¬ v.IsSharp
\end{lstlisting}

\noindent Definition~\ref{def:tensor-valuation}.
\phantomsection\label{lean:tensor-valuation}
\begin{lstlisting}
def IsTensorValuation (v : Ω → ℝ≥0∞) (t : Ω → Ω → Ω) : Prop :=
  ∀ a b, v a + v b = v (a ⊔ b) + v (t a b)
\end{lstlisting}

\noindent Theorem~\ref{thm:tensor-collapse}.
\phantomsection\label{lean:tensor-collapse}
\begin{lstlisting}
theorem IsTensorValuation.collapse {v : Ω → ℝ≥0∞} {t : Ω → Ω → Ω}
    (h : IsTensorValuation v t) {a : Ω} (ha : v a ≠ ⊤) :
    v (t a a) = v a
\end{lstlisting}

\noindent Example~\ref{ex:tensor-witness}.
\phantomsection\label{lean:tensor-witness}
\begin{lstlisting}
theorem forced_constant_of_add_valuation {v : ℕ → ℝ≥0∞}
    (h : IsTensorValuation v (· + ·)) (hfin : ∀ n, v n ≠ ⊤) (n : ℕ) :
    v n = v 0
\end{lstlisting}

\noindent Corollary~\ref{cor:tensor-dichotomy}.
\phantomsection\label{lean:tensor-dichotomy}
\begin{lstlisting}
theorem not_isTensorValuation_of_discriminates {v : Ω → ℝ≥0∞}
    {t : Ω → Ω → Ω} {a : Ω} (ha : v a ≠ ⊤) (hne : v (t a a) ≠ v a) :
    ¬ IsTensorValuation v t
\end{lstlisting}
\subsection{Probability Calculi from Certain States}

\phantomsection\label{lean:mix}
\begin{lstlisting}
def IsSharpVec (x : E → ℝ≥0∞) : Prop := ∀ e, x e = 0 ∨ x e = 1

def Mix (S : Set (E → ℝ≥0∞)) : Set (E → ℝ≥0∞) :=
  {x | ∃ (n : ℕ) (w : Fin n → ℝ≥0∞) (s : Fin n → E → ℝ≥0∞),
    (∀ i, s i ∈ S) ∧ ∑ i, w i = 1 ∧ x = ∑ i, w i • s i}
\end{lstlisting}

\phantomsection\label{lean:sharp-mem-mix-iff}
\begin{lstlisting}
theorem sharp_mem_Mix_iff {S : Set (E → ℝ≥0∞)} (hS : ∀ s ∈ S, IsSharpVec s)
    {x : E → ℝ≥0∞} (hx : IsSharpVec x) : x ∈ Mix S ↔ x ∈ S
\end{lstlisting}

\phantomsection\label{lean:mix-mix}
\begin{lstlisting}
theorem Mix_Mix (S : Set (E → ℝ≥0∞)) : Mix (Mix S) = Mix S
\end{lstlisting}

\phantomsection\label{lean:affinelaw}
\begin{lstlisting}


structure AffineLaw (E : Type*) [Fintype E] where
  a : E → ℝ≥0∞
  b : E → ℝ≥0∞
  k : ℝ≥0∞
  k' : ℝ≥0∞

def AffineLaw.Holds [Fintype E] (L : AffineLaw E) (x : E → ℝ≥0∞) : Prop :=
  ∑ e, L.a e * x e + L.k ≤ ∑ e, L.b e * x e + L.k'
\end{lstlisting}

\phantomsection\label{lean:affinelaw-holds-of-mix}
\begin{lstlisting}
theorem AffineLaw.holds_of_Mix [Fintype E] (L : AffineLaw E) {S : Set (E → ℝ≥0∞)}
    (hL : ∀ s ∈ S, L.Holds s) {x : E → ℝ≥0∞} (hx : x ∈ Mix S) : L.Holds x
\end{lstlisting}

\phantomsection\label{lean:mix-eq-laws}
\begin{lstlisting}
theorem Mix_eq_laws' {S : Set (E → ℝ≥0∞)} (hS : ∀ s ∈ S, IsSharpVec s) :
    Mix S = {x | ∀ L : AffineLaw E, (∀ s ∈ S, L.Holds s) → L.Holds x}
\end{lstlisting}

\phantomsection\label{lean:frame-recipe}
\begin{lstlisting}
theorem frame_recipe :
    Set.range (fun v : Valuation (LowerSet P) => v.toFun)
      = Mix {f | ∃ v : Valuation (LowerSet P), v.IsSharp ∧ f = v.toFun}
\end{lstlisting}

\phantomsection\label{lean:eq-deltapoint-of-issharp}
\begin{lstlisting}
theorem Valuation.eq_deltaPoint_of_isSharp {v : Valuation (LowerSet P)} (hv : v.IsSharp) :
    ∃ p, v = deltaPoint p
\end{lstlisting}

\phantomsection\label{lean:laws-iff-local-mix}
\begin{lstlisting}
structure FinLaw (E : Type*) where
  F : Finset E
  a : E → ℝ≥0∞
  b : E → ℝ≥0∞
  k : ℝ≥0∞
  k' : ℝ≥0∞

def FinLaw.Holds (L : FinLaw E) (x : E → ℝ≥0∞) : Prop :=
  ∑ e ∈ L.F, L.a e * x e + L.k ≤ ∑ e ∈ L.F, L.b e * x e + L.k'

def restr (F : Finset E) (x : E → ℝ≥0∞) : F → ℝ≥0∞ := fun e => x e

theorem laws_iff_local_Mix {S : Set (E → ℝ≥0∞)} (hS : ∀ s ∈ S, IsSharpVec s)
    (x : E → ℝ≥0∞) :
    (∀ L : FinLaw E, (∀ s ∈ S, L.Holds s) → L.Holds x) ↔
      ∀ F : Finset E, restr F x ∈ Mix (restr F '' S)
\end{lstlisting}

\phantomsection\label{lean:laws-iff-barycenter}
\begin{lstlisting}
def Cfg (Ev : Type) := Ev → Bool

def Good (S : Set (Ev → ℝ≥0∞)) (F : Finset Ev) (σ : Cfg Ev) : Prop :=
  ∃ s ∈ S, ∀ e : F, s e.1 = vec σ e.1

def atom (e : Ev) : Set (Cfg Ev) := {σ | σ e = true}

def bad (S : Set (Ev → ℝ≥0∞)) (F : Finset Ev) : Set (Cfg Ev) := {σ | ¬ Good S F σ}

theorem laws_iff_barycenter {S : Set (Ev → ℝ≥0∞)} (hS : ∀ s ∈ S, IsSharpVec s)
    (x : Ev → ℝ≥0∞) :
    (∀ L : FinLaw Ev, (∀ s ∈ S, L.Holds s) → L.Holds x) ↔
      ∃ μ : ProbabilityMeasure (Cfg Ev), (∀ F, (μ : Measure (Cfg Ev)) (bad S F) = 0) ∧
        ∀ e, x e = (μ : Measure (Cfg Ev)) (atom e)
\end{lstlisting}

\phantomsection\label{lean:mem-lin-iff}
\begin{lstlisting}
def flat (X : Type*) [Fintype X] : Set (X → ℝ≥0∞) := {x | ∑ a, x a ≤ 1}

noncomputable def app (t : X × Y → ℝ≥0∞) (x : X → ℝ≥0∞) : Y → ℝ≥0∞ :=
  fun b => ∑ a, x a * t (a, b)

def lin (X Y : Type*) [Fintype X] [Fintype Y] : Set (X × Y → ℝ≥0∞) :=
  {t | ∀ x ∈ flat X, app t x ∈ flat Y}

theorem mem_lin_iff [DecidableEq X] (t : X × Y → ℝ≥0∞) :
    t ∈ lin X Y ↔ ∀ a, ∑ b, t (a, b) ≤ 1
\end{lstlisting}

\phantomsection\label{lean:ispcs-lin}
\begin{lstlisting}
def orth (S : Set (E → ℝ≥0∞)) : Set (E → ℝ≥0∞) := {y | ∀ x ∈ S, pairing x y ≤ 1}

structure IsPCS (P : Set (E → ℝ≥0∞)) : Prop where
  closed : orth (orth P) = P
  total : ∀ e, ∃ x ∈ P, x e ≠ 0
  bounded : ∀ e, ∃ M ≠ ⊤, ∀ x ∈ P, x e ≤ M

theorem isPCS_lin : IsPCS (lin X Y)

theorem isPCS_flat : IsPCS (flat X)

theorem isPCS_tensDual : IsPCS tensDual
\end{lstlisting}

\phantomsection\label{lean:pcs-recipe}
\begin{lstlisting}
theorem pcs_recipe {X Y : Type} [Fintype X] [Fintype Y] [DecidableEq X] :
    lin X Y = Mix (Set.range (det : (X → Option Y) → X × Y → ℝ≥0∞))
\end{lstlisting}

\phantomsection\label{lean:sharp-lin-iff-clique}
\begin{lstlisting}
theorem sharp_lin_iff_clique {X Y : Type*} [Fintype X] [Fintype Y] [DecidableEq X] [DecidableEq Y]
    {t : X × Y → ℝ≥0∞} (h01 : ∀ p, t p = 0 ∨ t p = 1) :
    t ∈ lin X Y ↔ IsClique (arrowC (flatC X) (flatC Y)) (t · = 1)
\end{lstlisting}

\phantomsection\label{lean:sharp-td-iff}
\begin{lstlisting}
def tD (X Y X' Y' : Type) [Fintype X] [Fintype Y] [Fintype X'] [Fintype Y'] :
    Set ((X × Y) × (X' × Y') → ℝ≥0∞) :=
  {t | ∀ y ∈ lin X Y, ∀ z ∈ lin X' Y', ∑ p, ∑ q, y p * t (p, q) * z q ≤ 1}

def FTree (a : X) (cf : Y → X') (pq : (X × Y) × (X' × Y')) : Prop :=
  pq.1.1 = a ∧ pq.2.1 = cf pq.1.2

theorem sharp_tD_iff [Nonempty X] [Nonempty X'] {z : (X × Y) × (X' × Y') → ℝ≥0∞}
    (h01 : ∀ pq, z pq = 0 ∨ z pq = 1) :
    z ∈ tD X Y X' Y' ↔
      (∃ a cf, ∀ pq, z pq = 1 → FTree a cf pq) ∨ (∃ c af, ∀ pq, z pq = 1 → GTree c af pq)
\end{lstlisting}

\phantomsection\label{lean:sharp-mem-sequentialg}
\begin{lstlisting}
theorem sharp_mem_sequentialG [Nonempty X] [Nonempty X'] {z : (X × Y) × (X' × Y') → ℝ≥0∞}
    (hz : z ∈ tD X Y X' Y') (h01 : ∀ u, z u = 0 ∨ z u = 1) : z ∈ sequentialG X Y X' Y'
\end{lstlisting}

\phantomsection\label{lean:recipe-eq-sequentialg}
\begin{lstlisting}
structure Core (A B C D : Type) [Fintype A] [Fintype C] where
  μ : A → ℝ≥0∞
  κ : A → B → C → ℝ≥0∞
  ρ : A → B → C → D → ℝ≥0∞
  hκ : ∀ a b, ∑ c, κ a b c ≤ 1
  hρ : ∀ a b c d, ρ a b c d ≤ 1

def sequentialG (X Y X' Y' : Type) [Fintype X] [Fintype X'] : Set ((X × Y) × (X' × Y') → ℝ≥0∞) :=
  {t | ∃ (F : Core X Y X' Y') (G : Core X' Y' X Y), F.mass + G.mass ≤ 1 ∧
    t = fun pq => F.val pq.1.1 pq.1.2 pq.2.1 pq.2.2 + G.val pq.2.1 pq.2.2 pq.1.1 pq.1.2}

theorem recipe_eq_sequentialG [Nonempty X] [Nonempty X'] :
    Mix {z | z ∈ tD X Y X' Y' ∧ ∀ u, z u = 0 ∨ z u = 1} = sequentialG X Y X' Y'
\end{lstlisting}

\phantomsection\label{lean:sequentialg-subset-mix}
\begin{lstlisting}
theorem sequentialG_subset_Mix : sequentialG X Y X' Y' ⊆ Mix (sharpSeqG X Y X' Y')
\end{lstlisting}

\phantomsection\label{lean:mix-sequentialg}
\begin{lstlisting}
theorem Mix_sequentialG : Mix (sequentialG X Y X' Y') = sequentialG X Y X' Y'
\end{lstlisting}

\phantomsection\label{lean:pente-not-mix}
\begin{lstlisting}
theorem pentE_not_mix (hx : Function.Injective ex) (hy : Function.Injective ey)
    (hx' : Function.Injective ex') (hy' : Function.Injective ey')
    {n : ℕ} (w : Fin n → ℝ≥0∞) (hw : ∑ i, w i = 1) (z : Fin n → (X × Y) × (X' × Y') → ℝ≥0∞)
    (hz : ∀ i, z i ∈ tD X Y X' Y') (h01 : ∀ i u, z i u = 0 ∨ z i u = 1) :
    pentE ex ey ex' ey' ≠ ∑ i, w i • z i

theorem pentE_mem (hx : Function.Injective ex) (hy : Function.Injective ey)
    (hx' : Function.Injective ex') (hy' : Function.Injective ey') :
    pentE ex ey ex' ey' ∈ tD X Y X' Y'
\end{lstlisting}

\phantomsection\label{lean:pent-cycle}
\begin{lstlisting}
theorem pent_clique_le_two : ∀ B ∈ pent.powerset,
    (∀ s ∈ B, ∀ t ∈ B, compat2 s t) → B.card ≤ 2

theorem pent_indep_le_two : ∀ B ∈ pent.powerset,
    (∀ s ∈ B, ∀ t ∈ B, compat2 s t → s = t) → B.card ≤ 2
\end{lstlisting}

\phantomsection\label{lean:pent-bounds}
\begin{lstlisting}
def tensDual : Set (W × W → ℝ≥0∞) :=
  {t | ∀ y ∈ lin Bool Bool, ∀ z ∈ lin Bool Bool, ∑ p, ∑ q, y p * t (p, q) * z q ≤ 1}

def sequential : Set (W × W → ℝ≥0∞) :=
  {t | ∃ (α β : ℝ≥0∞) (F : FFirst) (G : GFirst), α + β ≤ 1 ∧ t = α • F.den + β • G.den}

theorem seq_pent_le_two {t : W × W → ℝ≥0∞} (ht : t ∈ sequential) : ∑ pq ∈ pent, t pq ≤ 2

theorem seq_pent_tight : ∃ t ∈ sequential, ∑ pq ∈ pent, t pq = 2

theorem pentElt_pent_sum : ∑ pq ∈ pent, pentElt pq = 5 * 2⁻¹

theorem tensDual_pent_le {t : W × W → ℝ≥0∞} (ht : t ∈ tensDual) :
    2 * ∑ pq ∈ pent, t pq ≤ 5
\end{lstlisting}

\phantomsection\label{lean:sharp-bang-mixed-zero}
\begin{lstlisting}
theorem sharp_bang_mixed_zero {z : ℕ × ℕ → ℝ≥0∞} (hz : z ∈ bang) (h01 : ∀ m, z m = 0 ∨ z m = 1)
    {i j : ℕ} (hi : 1 ≤ i) (hj : 1 ≤ j) : z (i, j) = 0
\end{lstlisting}

\phantomsection\label{lean:prom-half-not-mix}
\begin{lstlisting}
noncomputable def prom (x : Bool → ℝ≥0∞) : ℕ × ℕ → ℝ≥0∞ := fun m => x true ^ m.1 * x false ^ m.2

def bang : Set (ℕ × ℕ → ℝ≥0∞) := orthI (orthI gens)

theorem prom_half_not_mix : prom half ∈ bang ∧
    prom half ∉ Mix {z | z ∈ bang ∧ ∀ m, z m = 0 ∨ z m = 1}
\end{lstlisting}

\phantomsection\label{lean:td-eq-qstab}
\begin{lstlisting}
def excl (s t : (X × Y) × (X' × Y')) : Prop := s ≠ t ∧ cp2 s t

theorem tD_eq_QSTAB : tD X Y X' Y' = QSTAB (excl (X := X) (Y := Y) (X' := X') (Y' := Y'))
\end{lstlisting}

\phantomsection\label{lean:recipe-eq-stab}
\begin{lstlisting}
theorem recipe_eq_STAB :
    Mix {z | z ∈ tD X Y X' Y' ∧ ∀ u, z u = 0 ∨ z u = 1}
      = STAB (excl (X := X) (Y := Y) (X' := X') (Y' := Y'))
\end{lstlisting}

\phantomsection\label{lean:berge-of-stab-eq-qstab}
\begin{lstlisting}
def QSTAB : Set (V → ℝ≥0∞) := {x | ∀ C, IsCliqueG adj C → ∑ v ∈ C, x v ≤ 1}

def STAB : Set (V → ℝ≥0∞) := Mix {x | ∃ I, IsStableG adj I ∧ x = ind I}

structure Hole (adj : V → V → Prop) (n : ℕ) where
  h : ZMod n → V
  inj : Function.Injective h
  adj_iff : ∀ i j, i ≠ j → (adj (h i) (h j) ↔ cyc i j)

theorem berge_of_stab_eq_qstab (heq : STAB adj = QSTAB adj) (k : ℕ) (hk : 2 ≤ k)
    [NeZero (2 * k + 1)] : IsEmpty (Hole adj (2 * k + 1)) ∧ IsEmpty (Antihole adj (2 * k + 1))
\end{lstlisting}

\phantomsection\label{lean:berge-of-recipe-eq-pcs}
\begin{lstlisting}
theorem berge_of_recipe_eq_pcs
    (heq : Mix {z | z ∈ tD X Y X' Y' ∧ ∀ u, z u = 0 ∨ z u = 1} = tD X Y X' Y')
    (k : ℕ) (hk : 2 ≤ k) [NeZero (2 * k + 1)] :
    IsEmpty (Hole (excl (X := X) (Y := Y) (X' := X') (Y' := Y')) (2 * k + 1)) ∧
      IsEmpty (Antihole (excl (X := X) (Y := Y) (X' := X') (Y' := Y')) (2 * k + 1))
\end{lstlisting}

\subsection{Multiplicative--Additive Formulas}

\phantomsection\label{lean:linear-fm}
\begin{lstlisting}
inductive Fm where
  | base (n : ℕ)
  | neg (A : Fm)
  | tens (A B : Fm)
  | wth (A B : Fm)

def web : Fm → Type
  | base n => Fin n
  | neg A => web A
  | tens A B => web A × web B
  | wth A B => web A ⊕ web B

def P : (A : Fm) → Set (web A → ℝ≥0∞)
  | base n => flat (Fin n)
  | neg A => orth (P A)
  | tens A B => orth (orth {z | ∃ x ∈ P A, ∃ y ∈ P B, z = tvec x y})
  | wth A B => {z | (fun a => z (Sum.inl a)) ∈ P A ∧ (fun b => z (Sum.inr b)) ∈ P B}

def coh : (A : Fm) → web A → web A → Prop
  | base _ => fun x y => x = y
  | neg A => fun x y => x = y ∨ ¬ coh A x y
  | tens A B => fun s t => coh A s.1 t.1 ∧ coh B s.2 t.2
  | wth A B => fun s t => match s, t with
    | Sum.inl a, Sum.inl a' => coh A a a'
    | Sum.inr b, Sum.inr b' => coh B b b'
    | _, _ => True
\end{lstlisting}

\phantomsection\label{lean:ispcs-p}
\begin{lstlisting}
theorem isPCS_P : ∀ A : Fm, IsPCS (P A)
\end{lstlisting}

\phantomsection\label{lean:linear-invariants}
\begin{lstlisting}
theorem invariants : ∀ A : Fm,
    (∀ s : web A, Pi.single s 1 ∈ P A ∧ ∀ x ∈ P A, x s ≤ 1) ∧
    (∀ s t : web A, s ≠ t → coh A s t → pr s t ∈ P A) ∧
    (∀ s t : web A, s ≠ t → ¬ coh A s t → ∀ x ∈ P A, x s + x t ≤ 1)

theorem single_mem_P (A : Fm) (s : web A) : Pi.single s 1 ∈ P A

theorem pair_mem_P (A : Fm) {s t : web A} (hst : s ≠ t) (h : coh A s t) : pr s t ∈ P A
\end{lstlisting}

\phantomsection\label{lean:sharp-isclique}
\begin{lstlisting}
theorem sharp_isClique (A : Fm) {z : web A → ℝ≥0∞} (hz : z ∈ P A)
    (h01 : ∀ s, z s = 0 ∨ z s = 1) : ∀ s t, z s = 1 → z t = 1 → coh A s t
\end{lstlisting}

\phantomsection\label{lean:recipe-subset-p}
\begin{lstlisting}
theorem recipe_subset_P (A : Fm) :
    Mix {z | z ∈ P A ∧ ∀ s, z s = 0 ∨ z s = 1} ⊆ P A

theorem recipe_sharp_iff (A : Fm) {x : web A → ℝ≥0∞} (hx : ∀ s, x s = 0 ∨ x s = 1) :
    x ∈ Mix {z | z ∈ P A ∧ ∀ s, z s = 0 ∨ z s = 1} ↔ x ∈ P A
\end{lstlisting}

\phantomsection\label{lean:mix-subset-orth}
\begin{lstlisting}
theorem Mix_subset_orth {S T : Set (E → ℝ≥0∞)} (h : S ⊆ orth T) : Mix S ⊆ orth T
\end{lstlisting}

\phantomsection\label{lean:p-lolli-base}
\begin{lstlisting}
theorem P_lolli_base (n m : ℕ) : P (lolli (base n) (base m)) = lin (Fin n) (Fin m)
\end{lstlisting}

\phantomsection\label{lean:p-second}
\begin{lstlisting}
def second (n m n' m' : ℕ) : Fm := neg (tens (lolli (base n) (base m)) (lolli (base n') (base m')))

theorem P_second (n m n' m' : ℕ) :
    P (second n m n' m') = tD (Fin n) (Fin m) (Fin n') (Fin m')
\end{lstlisting}

\phantomsection\label{lean:recipe-formula}
\begin{lstlisting}
theorem recipe_first_order (n m : ℕ) :
    P (lolli (base n) (base m)) = ConstructiveProb.Mix (Set.range
      (det : (Fin n → Option (Fin m)) → Fin n × Fin m → ℝ≥0∞))

theorem recipe_second_order {n m n' m' : ℕ} (hn : 0 < n) (hn' : 0 < n') :
    ConstructiveProb.Mix {z | z ∈ P (second n m n' m') ∧ ∀ u, z u = 0 ∨ z u = 1}
      = sequentialG (Fin n) (Fin m) (Fin n') (Fin m')

theorem recipe_second_order_ssubset {n m n' m' : ℕ} (hn : 2 ≤ n) (hm : 2 ≤ m) (hn' : 2 ≤ n')
    (hm' : 2 ≤ m') :
    ConstructiveProb.Mix {z | z ∈ P (second n m n' m') ∧ ∀ u, z u = 0 ∨ z u = 1}
      ⊂ P (second n m n' m')
\end{lstlisting}

\phantomsection\label{lean:shannon-clique}
\begin{lstlisting}
abbrev S2 : Fm := second 2 2 2 2

abbrev TT : Fm := neg (tens S2 S2)

noncomputable def shannonF : Finset (web (tens S2 S2)) := shannon.image fun u => (φ u.1, φ u.2)

theorem shannonF_clique : ∀ u ∈ shannonF, ∀ v ∈ shannonF, coh TT u v
\end{lstlisting}

\phantomsection\label{lean:shannon-not-mem}
\begin{lstlisting}
theorem shannonF_not_mem : (fun u => if u ∈ shannonF then (1 : ℝ≥0∞) else 0) ∉ P TT
\end{lstlisting}

\subsection{Quantum Event Logic}

\phantomsection\label{lean:cabello-set}
\begin{lstlisting}
def vec : Fin 18 → Fin 4 → ℤ := ![
  ![-1, 1, 1, 1], ![0, 0, 0, 1], ![0, 0, 1, 0], ![0, 0, 1, 1], ![0, 1, -1, 0], ![0, 1, 0, -1],
  ![0, 1, 0, 0], ![1, -1, -1, 1], ![1, -1, 0, 0], ![1, -1, 1, -1], ![1, 0, -1, 0], ![1, 0, 0, -1],
  ![1, 0, 0, 1], ![1, 0, 1, 0], ![1, 1, -1, 1], ![1, 1, 0, 0], ![1, 1, 1, -1], ![1, 1, 1, 1]]

def bas : Fin 9 → Fin 4 → Fin 18 := ![
  ![1, 2, 15, 8], ![1, 6, 13, 10], ![9, 7, 15, 3], ![9, 17, 10, 5], ![2, 6, 12, 11],
  ![7, 17, 11, 4], ![14, 16, 8, 3], ![14, 0, 13, 5], ![16, 0, 12, 4]]

def dot (u v : Fin 4 → ℤ) : ℤ := ∑ k, u k * v k
\end{lstlisting}

\phantomsection\label{lean:basis-orth}
\begin{lstlisting}
theorem basis_orth : ∀ b : Fin 9, ∀ i j : Fin 4, i ≠ j → dot (vec (bas b i)) (vec (bas b j)) = 0

theorem occ_two : ∀ v : Fin 18,
    (univ.filter (fun p : Fin 9 × Fin 4 => bas p.1 p.2 = v)).card = 2
\end{lstlisting}

\phantomsection\label{lean:no-ks}
\begin{lstlisting}
theorem no_ks : ¬ ∃ f : Fin 18 → Bool, ∀ b, (univ.filter (fun i => f (bas b i))).card = 1
\end{lstlisting}

\phantomsection\label{lean:no-certain-state}
\begin{lstlisting}
def certain : Set (Fin 18 → ℝ≥0∞) :=
  {s | IsSharpVec s ∧ ∀ b, ∑ i, s (bas b i) = 1}

theorem no_certain_state : certain = ∅

theorem Mix_certain_empty : Mix certain = ∅
\end{lstlisting}

\phantomsection\label{lean:parseval}
\begin{lstlisting}
theorem parseval (ψ : Fin 4 → ℝ) (b : Fin 9) :
    ∑ i, (∑ k, ψ k * vec (bas b i) k) ^ 2 / ∑ k, ((vec (bas b i) k : ℝ)) ^ 2 = ∑ k, ψ k ^ 2
\end{lstlisting}

\phantomsection\label{lean:born-not-mix}
\begin{lstlisting}
def states : Set (Fin 18 → ℝ≥0∞) := {x | ∀ b, ∑ i, x (bas b i) = 1}

noncomputable def born (ψ : Fin 4 → ℝ) : Fin 18 → ℝ≥0∞ := fun v =>
  ENNReal.ofReal ((∑ k, ψ k * vec v k) ^ 2 / ∑ k, ((vec v k : ℝ)) ^ 2)

theorem born_mem {ψ : Fin 4 → ℝ} (hψ : ∑ k, ψ k ^ 2 = 1) : born ψ ∈ states

theorem born_not_mix {ψ : Fin 4 → ℝ} (hψ : ∑ k, ψ k ^ 2 = 1) :
    born ψ ∈ states ∧ born ψ ∉ Mix certain
\end{lstlisting}

\subsection{Finite Frames without Choice}

\phantomsection\label{lean:exists-sharp-separating-fin}
\begin{lstlisting}
def JoinPrime (j : α) : Prop := ∀ x y : α, j ≤ x ⊔ y → j ≤ x ∨ j ≤ y

theorem exists_joinPrime_aux : ∀ (n : ℕ) (a b : α), (below a).card = n → ¬ a ≤ b →
    ∃ j, j ≤ a ∧ ¬ j ≤ b ∧ JoinPrime j

theorem exists_sharp_separating_fin [BoundedOrder α] {a b : α} (hab : ¬ a ≤ b) :
    ∃ v : α → ℕ, (∀ x, v x = 0 ∨ v x = 1) ∧ Monotone v ∧ v ⊥ = 0 ∧ v ⊤ = 1 ∧
      (∀ x y, v x + v y = v (x ⊔ y) + v (x ⊓ y)) ∧ v a = 1 ∧ v b = 0
\end{lstlisting}

\phantomsection\label{lean:em-of-sharp-complement-fin}
\begin{lstlisting}
theorem em_of_sharp_complement_fin [BoundedOrder α] (a c : α)
    (h : ∀ v : α → ℕ, (∀ x, v x = 0 ∨ v x = 1) → Monotone v → v ⊥ = 0 → v ⊤ = 1 →
      (∀ x y, v x + v y = v (x ⊔ y) + v (x ⊓ y)) → v a + v c = 1) :
    a ⊔ c = ⊤
\end{lstlisting}

\subsection{Choice-Free Finite Results}

\phantomsection\label{lean:constr-sierpinski}
\begin{lstlisting}
inductive Op | bot | halts | top
  deriving DecidableEq

def neg : Op → Op
  | bot => top
  | halts => bot
  | top => bot

theorem neg_spec : ∀ o o' : Op, (inf o' o = bot) ↔ le o' (neg o) = true

theorem neg_halts : neg halts = bot

theorem frontier_iff : ∀ x : Bool, frontier x = true ↔ x = false

theorem slack_eq (wT wF : Nat) :
    val wT wF top = val wT wF halts + val wT wF (neg halts) + wt (frontier false) wF
\end{lstlisting}

\phantomsection\label{lean:constr-boundary}
\begin{lstlisting}
def ext (ops : List (Fin n → Bool)) (U : Fin n → Bool) (x : Fin n) : Bool :=
  ops.any fun O => O x && allFin n (fun y => !(O y && U y))

def bd (ops : List (Fin n → Bool)) (U : Fin n → Bool) (x : Fin n) : Bool :=
  !U x && !ext ops U x

theorem sub_ext {ops : List (Fin n → Bool)} {U : Fin n → Bool} {O : Fin n → Bool}
    (hO : O ∈ ops) (hd : ∀ y, (!(O y && U y)) = true) {x : Fin n} (hx : O x = true) :
    ext ops U x = true

theorem slack_eq_boundary_fin (ops : List (Fin n → Bool)) (U : Fin n → Bool) (μ : Fin n → Nat) :
    fs n μ = fs n (fun x => wt (U x) (μ x)) + fs n (fun x => wt (ext ops U x) (μ x)) +
      fs n (fun x => wt (bd ops U x) (μ x))
\end{lstlisting}

\phantomsection\label{lean:constr-kraft}
\begin{lstlisting}
def ksum (L : Nat) : List (List Bool) → Nat
  | [] => 0
  | w :: S => 2 ^ (L - w.length) + ksum L S

theorem kraft : ∀ (L : Nat) (S : List (List Bool)), PF S → (∀ w ∈ S, w.length ≤ L) →
    ksum L S ≤ 2 ^ L
\end{lstlisting}

\phantomsection\label{lean:constr-farkas}
\begin{lstlisting}
structure Row (n : Nat) where
  c : Fin n → Int
  s : Bool

inductive Der {n : Nat} (rows : List (Row n)) : Row n → Prop
  | base {r : Row n} : r ∈ rows → Der rows r
  | comb {r₁ r₂ : Row n} (p q : Int) : 0 ≤ p → 0 ≤ q → Der rows r₁ → Der rows r₂ →
      Der rows (comb p q r₁ r₂)

theorem fm : ∀ (n : Nat) (rows : List (Row n)), (∃ w, ∀ r ∈ rows, Sat r w) ∨ Der rows (zeroS n)

theorem complete (hd : 0 < d)
    (hlaw : ∀ (c : Fin m → Int) (k : Int), (∀ j, k ≤ dot m c (S j)) → k * d ≤ dot m c y) :
    ∃ (lam : Fin N → Int) (t : Int), 0 < t ∧ (∀ j, 0 ≤ lam j) ∧
      dot N (fun _ => 1) lam = t ∧ ∀ e, d * dot N lam (fun j => S j e) = t * y e

theorem sound (hd : 0 < d) (lam : Fin N → Int) (t : Int) (ht : 0 < t) (hl : ∀ j, 0 ≤ lam j)
    (hsum : dot N (fun _ => 1) lam = t) (hx : ∀ e, d * dot N lam (fun j => S j e) = t * y e)
    (c : Fin m → Int) (k : Int) (hc : ∀ j, k ≤ dot m c (S j)) : k * d ≤ dot m c y
\end{lstlisting}

\phantomsection\label{lean:constr-fubini}
\begin{lstlisting}
structure EOD (A : Nat → Nat → Bool) where
  N : Nat
  lt : Bool
  eq : Bool
  gt : Bool
  spec : ∀ m m', N ≤ m → N ≤ m' →
    A m m' = if m < m' then lt else if m = m' then eq else gt

def L {A : Nat → Nat → Bool} (w : EOD A) : Bool := w.lt
/-- The iterated limit with `m → ∞` first. -/
def R {A : Nat → Nat → Bool} (w : EOD A) : Bool := w.gt

def R {A : Nat → Nat → Bool} (w : EOD A) : Bool := w.gt

theorem L_modular {A B} (a : EOD A) (b : EOD B) :
    val (L a) + val (L b) = val (L (union a b)) + val (L (inter a b))

theorem rect_L {S T} (s : EC S) (t : EC T) : L (rect s t) = (s.v && t.v)

theorem rect_R {S T} (s : EC S) (t : EC T) : R (rect s t) = (s.v && t.v)

theorem products_differ : L tri ≠ R tri
\end{lstlisting}
